\section{Introduction}\label{sec:intro}

\subsection{Motivation}\label{sec:motivation}

When a neural network is reported to reach, say, $98\%$ accuracy on a benchmark,
that figure is an estimate from one finite test set. It carries no stated
confidence, it varies from one test split to another, and it can be quietly
inflated when the same test set is consulted repeatedly during model selection.
For a deployed classifier we often want something stronger: ``with probability
at least $99\%$, the expected error of this model is below $r$''. A bound of this
kind is a \emph{risk certificate}, computed from the training data and the
learning algorithm rather than from a held-out test set.

PAC-Bayes theory \parencite{mcallester1999,seeger2002,catoni2007} is the main
route to such guarantees, for \emph{stochastic} predictors whose weights are
drawn from a distribution $Q$ learnt from data.
For a long time these bounds were of theoretical interest only: applied to a
modern over-parameterised network they returned values above one, which certify
nothing. That changed when the bound was used not merely to analyse a
trained model but to \emph{train} it. PAC-Bayes-with-Backprop (PBB)
\parencite{perezortiz2021} optimises the weight distribution by minimising a
PAC-Bayes bound directly, obtaining certificates that are numerically meaningful
and competitive in accuracy. This project reproduces that framework,
extends it to two further datasets, and examines how trustworthy the resulting
certificates are once the computation behind them is scrutinised.

\subsection{Why certificate reproduction is hard}\label{sec:problem}

A test-error number is forgiving: a small mistake in computing it usually produces
a slightly wrong but still plausible figure. A certificate is not. Its validity
rests on a chain of conditions that must all hold at once. The empirical risk
inside the bound has to be measured on data the prior never saw, and the sample
size dividing the complexity term has to be the size of that held-out set rather
than of the whole dataset. The risk of a stochastic predictor has no closed form,
so it is estimated by Monte-Carlo sampling, which introduces a second failure
probability that must be combined with the PAC-Bayes one. And the object certified
is the stochastic (Gibbs) classifier, not the deterministic network a practitioner
deploys. Handle any of these wrongly and the pipeline still produces a number that
looks like a certificate, but it no longer holds.

We therefore re-implement the pipeline from first principles alongside
reproducing and extending PBB, and test each condition explicitly. That surfaced
several points at which the reference implementation departs from them; we
correct those and recompute every certificate accordingly.

\subsection{Subject area and key related work}\label{sec:related}

Non-vacuous numerical bounds go back to \textcite{langford2001}, for small
networks only. The modern line begins with \textcite{dziugaite2017}: a
deterministic network has a point-mass weight distribution, for which the
PAC-Bayes complexity term is infinite, so by spreading the weights into a Gaussian
and optimising the spread they turned an infinite penalty into a finite one and
reported the first non-vacuous bound (around $0.21$) on MNIST. It was loose, but
it established that the approach was numerically viable.

Two factors separate a loose bound from a tight one, and both matter here. The
first is keeping the posterior $Q$ close to the prior $P$, since the bound
penalises their KL divergence, and a prior fixed before seeing any data sits far
from any well-trained posterior. \textcite{dziugaite2018} addressed this by making
the prior depend on the data through differential privacy, and
\textcite{rivasplata2018} analysed instance-dependent priors more generally; both
show that moving the prior towards the data shrinks the KL substantially. The
second is to optimise the bound itself. \textcite{perezortiz2021}, the work this
project reproduces, trains $Q$ by minimising one of three surrogate PAC-Bayes
objectives (a McAllester/Pinsker form $f_{\text{classic}}$, the Catoni-style bound
of \textcite{thiemann2017} $f_{\lambda}$, and a quadratic relaxation
$f_{\text{quad}}$), combined with a prior learnt on a held-out portion of the
training data, and reports $0$--$1$ certificates around $0.02$--$0.03$ on MNIST,
an order of magnitude tighter than \textcite{dziugaite2017}, at good accuracy.
Later work develops the self-certified and learnt-prior settings
\parencite{perezortiz2021b,perezortiz2021c}, and \textcite{garciaperez2025}
tightens the underlying inequalities.

Which predictor the bound controls matters throughout what follows: it is a
statement about the Gibbs classifier, which resamples weights for each
prediction. Transferring it to the deterministic majority vote is possible in
principle through the C-bound \parencite{lacasse2006,germain2015}, which relates
the majority-vote risk to the mean and variance of the Gibbs risk, but that is a
different and generally looser statement. The baseline we compare against,
Bayes-by-Backprop \parencite{blundell2015}, learns a Gaussian posterior by
maximising a variational objective rather than a risk bound, and so achieves
strong accuracy without a tight certificate. More broadly, deep ensembles
\parencite{lakshminarayanan2017}, MC-dropout \parencite{gal2016} and post-hoc
calibration \parencite{guo2017} quantify uncertainty empirically without a
distribution-free guarantee on the risk, while randomised smoothing
\parencite{cohen2019} certifies a different, adversarial quantity.

\subsection{Objectives}\label{sec:rqs}

The project brief is to reproduce the PBB experiments on MNIST, extend them to
Fashion-MNIST, and analyse CIFAR-10. Within that brief the technical work pursues
four questions. The first is one of \emph{reproduction fidelity}: under an
explicitly reduced compute budget, how closely does a corrected, independent
re-implementation match the reference's MNIST results when the same metrics are
compared? To make that decidable rather than a matter of impression, targets and
tolerance are fixed in advance: four MNIST cells of the reference's Table~1 --- FCN
with a data-free prior under $f_{\text{quad}}$, FCN with a learnt prior under
$f_{\text{quad}}$ and under $f_{\lambda}$, and CNN with a learnt prior under
$f_{\text{quad}}$ --- each counting as reproduced when our mean stochastic test
error is within one percentage
point of the published value. We set \emph{no} tolerance on the certificate,
because one computed at $m=2{,}000$ is not the same quantity as one computed at
$m=150{,}000$; instead we report the gap per cell and measure how much of it the
Monte-Carlo budget explains. The second concerns the \emph{anatomy} of the certificate, namely how
the empirical Gibbs risk, the KL term and the bound-set size $n_{\mathrm b}$ each
contribute to the final value across the three datasets, and how a learnt prior
shifts that balance. Building on this, the third asks how \emph{sensitive} the
certificate is to the choices around it: the prior family, the amount of training
data, the Monte-Carlo budget, the random seed, and the sample-count convention
that the audit corrects. The last isolates a single notion of difficulty: with the
architecture, sample size and compute held fixed, does injected label noise move
the empirical-risk, KL and certificate terms together, so that a difficulty effect
can be separated from the confounds of a cross-dataset comparison?

The dissertation makes five contributions. It (i) reproduces the principal PBB
MNIST cells under a documented, reduced-compute protocol and compares them with the
published numbers on identically-defined metrics; (ii) extends the same pipeline to
Fashion-MNIST and CIFAR-10, adding the learnt-versus-random prior contrast on every
dataset, a CIFAR-10 random-prior control the original study did not report, and
its 13- and 15-layer CIFAR-10 architectures, at the deeper of which the fixed
configuration begins to fail to train;
(iii) identifies, in the public reference \emph{implementation} and not in the
theory of \textcite{perezortiz2021}, whose inequalities we use unchanged, five
issues in the certificate pipeline, and corrects each under a unit test: three
change the certificate value (the bound-set sample count, the batched Monte-Carlo
aggregation, the joint confidence accounting), one affects reproducibility (seed
control) and one affects whether the learnt-prior construction meets the
independence condition at all (a label-independent partition). It also pays for a
step of its own, since the prior scale was chosen by comparing certificates on
the bound set, with a union over a grid of candidate prior scales, and states
precisely what that correction does and does not settle;
(iv) decomposes certificate behaviour into its empirical-risk, KL, sample-size and
Monte-Carlo components, by class, and under a controlled label-noise perturbation;
and (v) releases a reproducible artefact in which every results figure and
reported number is generated from a results registry with per-run provenance and an
inclusion/exclusion record.

\subsection{Report structure}\label{sec:structure}

Section~\ref{sec:theory} sets out the theory and the assumptions it rests on,
Section~\ref{sec:method} the experimental design and the corrected pipeline, and
Section~\ref{sec:results} the reproduction, the certificate anatomy, the
sensitivity analyses and the controlled-difficulty study.
Section~\ref{sec:evaluation} reflects on validity and limitations, and
Section~\ref{sec:conclusion} concludes.
